\شروع‌مسئله[]
(رزرو باند فرودگاه ۲)
مسئله رزرو باند فرودگاه که در کلاس مطرح شد را به یاد بیاورید. اوضاع برج مراقبت به هم ریخته و پروازها روی هم رزرو شده‌اند. حال از شما کمک می‌خواهند که یک داده‌ساختار در اختیار آن‌ها قرار دهید که با کمک آن بتوانند اوضاع را به حالت عادی برگردانند. این داده‌ساختار باید قادر باشد سه کار زیر را انجام دهد:
\شروع{فقرات}[label= - \ ]
\فقره
اضافه کردن یک بازه به داده‌ساختار شامل شماره پرواز، ابتدای بازه و انتهای بازه.
\فقره
حذف کردن بازه مربوط به یک پرواز از داده‌ساختار با گرفتن شماره پرواز.
\فقره
بررسی این که آیا یک بازه جدید با بازه‌های قبلی همپوشانی دارد یا خیر (اشتراک در نقطه ابتدا و انتها همپوشانی محسوب نمی‌شود).
\پایان{فقرات}

یک داده‌ساختار مناسب برای این کار پیشنهاد دهید.

\شروع‌مسئله[]
(زوج مؤدب)
در یک مجموعۀ $S$ از اعداد متمایز، به دنبال دو عضو مختلف هستیم که جمع آن‌ها در یک محدودۀ داده شده باشد. یعنی در ورودی اعضای مجموعۀ $S$ و دو مقدار $\ell$ و $r$ داده شده است و یک زوج از اعضای $S$ مؤدب هستند اگر و تنها اگر
$$\ell \leq s_i + s_j \leq r.$$
همچنین مجموعۀ $S$ این خاصیت را دارد که
$$\ln\Big(\max(S) - \min(S)\Big) < 40254.$$

الگوریتم بهینه‌ای برای پیدا کردن یک زوج مؤدب ارائه دهید.

\شروع‌مسئله[]
(شیروانی‌ها)
در این سؤال با یک دنیای دو بعدی سر و کار داریم که با صفحه دکارتی مدل شده است و زمین آن خط $y=0$ یا همان محور $x$ها است. در این صفحه $n$ سقف شیب‌دار (شیروانی) داریم که هر کدام معادل با یک پاره‌خط بالای محور $x$ها هستند. پاره‌خط مربوط به هیچ‌کدام از سقف‌ها عمودی و افقی نیست. همچنین سقف‌ها همدیگر را قطع نمی‌کنند. در این صفحه باران از بالا و به صورت عمودی شروع به بارش می‌کند. باران به صورت یکنواخت است یعنی یک واحد باران در هر واحد طول افقی در هر ثانیه. وقتی باران روی یک سقف می‌بارد، روی آن سقف حرکت کرده و از لبه پایین‌تر سقف سرریز می‌کند. با فرض این‌که مدتی از شروع بارش می‌گذرد و حالت سیستم در بالای محور $x$ها به تعادل رسیده است، می‌خواهیم ببینیم در هر نقطه از زمین (محور $x$ها) چقدر باران می‌ریزد.

ورودی مسئله شامل مختصات دو سر پاره‌خط‌های متناظر با شیروانی‌ها است و خروجی مسئله به این صورت است: محور $x$ برحسب میزان بارش به تعدادی بازه (ماکسیمال) بسته یا باز افراز می‌شود که در نقاط هر بازه، بارش یکسانی وجود دارد. بازه‌ها دو نوع هستند: بازه‌های با طول مثبت و بازه‌های با طول صفر که یک تک نقطه هستند. برای هر بازه که طول مثبت دارد میزان بارش بر واحد طول در هر ثانیه را خروجی می‌دهیم و برای هر بازه با طول صفر (یعنی یک نقطه) که بارش سرریز شده از یک سقف روز آن می‌ریز‌د، میزان بارش در هر ثانیه را خروجی می‌دهیم. بازه‌ها نیز باید به صورت مرتب شده بر اساس نقطه سمت چپ (و در حالت مساوی بودن نقطه‌های سمت چپ، بر اساس نقطه سمت راست) در خروجی آورده شوند. به طور مثال اگر دو شیروانی داشته باشیم که یکی بین دو نقطه $(0,10)$ و $(5,5)$ باشد و دیگری از $(4,10)$ تا $(11,11)$ باشد، باران به این صورت توزیع می‌شود:
\شروع{فقرات}[label= - \ ]
\فقره
در بازه $(-\infty,0)$ باران یک واحد در واحد طول در هر ثانیه می‌بارد.
\فقره
در بازه $[0, 5)$ بارانی نمی‌بارد.
\فقره
در بازه $[5, 5]$ باران ۱۱ واحد در ثانیه می‌بارد.
\فقره
در بازه $(5, 11]$ بارانی نمی‌بارد.
\فقره
در بازه $(11, \infty)$ باران یک واحد در واحد طول در هر ثانیه می‌بارد.
\پایان{فقرات}


 یک الگوریتم بهینه برای این مسئله طراحی کنید.

